\section{Design Overview}

We present an LLM inference system designed for mobile NPUs and optimized for test-time scaling workloads.

To address the accuracy challenge, we adopt 4-bit fine-grained group quantization for the primary weights while keeping activations in floating-point. During runtime, we dynamically dequantize the weights into floating-point values on the fly, leveraging the powerful FP16 matrix computation capabilities of the NPU to efficiently support test-time scaling tasks.

For the unavoidable general-purpose computations — where the vector processing units exhibit limited memory bandwidth and compute throughput, our core strategy includes:

\begin{itemize}
    \item Employing hardware-aware offline design to minimize runtime computation overhead;
    \item Fully exploiting the intrinsic capabilities of SIMD vector units to bridge the gap between specialized hardware and flexible software requirements.
\end{itemize}

Specifically, we introduce the following techniques:

\paragraph{Hardware-aware Fine-grained Tile Quantization Scheme.} We propose a novel quantization layout that performs group quantization in fine-grained rectangular tiles, as opposed to conventional approaches that group along the accumulation axis. To align with the memory access patterns of both the matrix and vector units, we introduce an offline pipeline involving weight pre-quantization transformation, quantization, and post-quantization transformation. This enhances the continuity of runtime memory access and eliminates unnecessary computational overhead.

\paragraph{Efficient LUT-Based Computation.} For more complex runtime operations, we leverage the vector unit’s lookup table (LUT) instructions and generalized LUT mechanisms to replace intricate transformation logic. This approach accelerates key bottleneck operations in test-time scaling workloads, including dequantization within mixed-precision GEMM and the Softmax operation in Attention.

% We propose an LLM inference system that targets Hexagon NPU and test-time scaling workloads. To tackle the hardware precision issue as well as achieving low memory overhead, we mainly adopt fine-grained 4-bit group quantization for weights. The model weights are dequantized to FP16 at runtime to be compatible with the FP16 HMX unit. Constrained by the relatively weak general computing capability of the HVX, we employ a two-stage approach. During the offline stage, we per-arrange the weights according to the data layout required by HMX before applying per-group post-training quantization to eliminate discontinuous memory accesses. A further layout permutation is applied to the model weights to facilitate wide vector registers. We also design efficient online operators. Generalized vector table lookup instructions are utilized to accelerate complex transformations in the time-consuming softmax and dequantization process.

% To achieve runtime computation scaling with low extra cost, we design HMX-centric GEMM (xx section) and Attention (xx section) operators. To meet the high-precision requirements of test-time scaling workloads (and the constraint of low memory usage \hzx{when should we mention this?}), we adopt fine-grained 4-bit group quantization for weights. The weights are dequantized to FP16 at runtime to be compatible with the floating-point HMX unit. Constrained by the relatively weak general computing capability of the HVX, we pre-arrange the weights offline according to the data layout required by HMX/HVX (xx section) to reduce the runtime dequantization overhead. We further employ efficient LUT-based dequantization (xx section) and a weight loading-dequantization pipeline (xx section) to accelerate the dequantization process. For the Attention operation, we keep most of the computations in FP16 and utilize a vector-gather-based softmax (xx section) to maintain both computational precision and speed.
